% F6 of article/paper-figures.md, redrawn, and the only drawing of the
% branching scheme.  One picture for the whole hierarchy: the three rungs,
% nested (level 2 inside the equality child of level 1, level 3 inside the
% Murty leaf); an integer line under the first split, so that "disjoint and
% exhaustive" is shown rather than asserted and the later rungs can be read
% off it; the construction as what happens at every node; and the rung that
% is deliberately not taken.  Greyscale only, straight edges only, every
% label clear of every edge, and the spine of the tree runs at x=7.5 -- the
% heading blocks are 66mm wide so that none of them reaches it.
\begin{figure}[!t]
\centering
\begin{tikzpicture}[>=Latex, x=1cm, y=1cm,
  nd/.style={rectangle, draw=black!75, fill=white, line width=0.5pt,
             inner sep=3pt, minimum width=26mm, minimum height=6mm,
             font=\footnotesize},
  fx/.style={nd, fill=black!12},
  rt/.style={nd, minimum width=52mm},
  rep/.style={rectangle, draw=black!60, densely dashed, fill=white,
              line width=0.5pt, inner sep=3pt, text width=40mm,
              align=flush center, font=\scriptsize},
  an/.style={font=\scriptsize, align=flush center, text width=60mm, anchor=north},
  hd/.style={font=\scriptsize, align=flush left, anchor=north west,
             text width=66mm},
  bnr/.style={rectangle, draw=black!45, fill=black!5, rounded corners=1pt,
              inner sep=4pt, text width=142mm, align=flush left, font=\scriptsize,
              anchor=north},
  bnx/.style={bnr, draw=black!60, densely dashed, fill=white},
  sp/.style={draw=black!75, line width=0.5pt},
  ed/.style={sp, ->},
  ax/.style={draw=black!55, line width=0.4pt},
  seg/.style={draw=black!85, line width=1.5pt},
  tk/.style={font=\tiny, inner sep=1pt}]

% ---------------------------------------------------------------- banner --
\node[bnr] (ban) at (7.5,1.62) {\textbf{At every node} the construction
  (Algorithm~\ref{alg:construct}) builds the plan that node is partitioned
  around, and its trucks join one shared pool: a column is valid in every
  region, so a pruned node leaves its columns behind.};

% =============================================================== level 1 ==
\node[hd] at (0.4,0.62) {\textbf{level 1} --- the total fleet
  $X=\sum_{g,t}n_{gt}$};

\node[rt] (root) at (7.5,-0.45)
  {root: the constructed plan $\mathbf{n}^{\star}$, $X^{\star}$ trucks};

\node[nd] (a1) at (3.0,-1.70) {$X\le X^{\star}-1$};
\node[fx] (b1) at (7.5,-1.70) {$X=X^{\star}$};
\node[nd] (c1) at (12.0,-1.70) {$X\ge X^{\star}+1$};
\coordinate (j1) at (7.5,-1.07);
\draw[sp] (root.south) -- (j1);
\draw[ed] (j1) -- (a1.north);
\draw[ed] (j1) -- (b1.north);
\draw[ed] (j1) -- (c1.north);

% the integer line: the three children abut at $X^\star$ and leave no gap
\draw[ax] (0.95,-2.55) -- (14.05,-2.55);
\foreach \i in {0,...,14} {\draw[ax] (1.2+\i*0.9,-2.55) -- ++(0,0.075);}
\draw[seg] (1.2,-2.55) -- (6.6,-2.55);
\draw[seg] (8.4,-2.55) -- (13.8,-2.55);
\fill (7.5,-2.55) circle (1.7pt);
\node[tk, anchor=east] at (1.05,-2.55) {$\cdots$};
\node[tk, anchor=west] at (13.95,-2.55) {$\cdots$};
\node[tk] at (6.6,-2.82) {$X^{\star}\!-\!1$};
\node[tk] at (8.4,-2.82) {$X^{\star}\!+\!1$};
\node[tk, anchor=west] at (7.66,-2.42) {$X^{\star}$};
\node[tk, anchor=east] at (13.8,-2.82) {$X\in\mathbb{Z}$};

\node[an] at (3.0,-3.07) {the only region that can hold a cheaper plan:
  popped first, and split again on its own total};
\node[an] at (12.0,-3.07) {a whole truck more, with nothing in the objective
  to offset it: usually pruned by bound};

% =============================================================== level 2 ==
\node[hd] at (0.4,-4.00) {\textbf{level 2} --- the fleet vector, inside
  $X=X^{\star}$: Murty at the first pair $j$ that may differ};

\node[nd] (a2) at (3.0,-5.45) {$n_{j}\le n^{\star}_{j}-1$};
\node[fx] (b2) at (7.5,-5.45) {$n_{j}=n^{\star}_{j}$};
\node[nd] (c2) at (12.0,-5.45) {$n_{j}\ge n^{\star}_{j}+1$};
\coordinate (j2) at (7.5,-4.85);
\draw[sp] (b1.south) -- (7.5,-2.55);
\draw[sp] (7.5,-2.55) -- (j2);
\draw[ed] (j2) -- (a2.north);
\draw[ed] (j2) -- (b2.north);
\draw[ed] (j2) -- (c2.north);

\node[an, text width=38mm] at (3.0,-5.95) {the same trichotomy, one pair at
  a time; empty when $n^{\star}_{j}=0$};
\node[an, text width=38mm] at (12.0,-5.95) {the total is fixed here, so a pair
  that moves up forces another down};

\node[rep] (rep) at (7.5,-6.37) {split again at the next pair,\\
  until every $n_{gt}$ is fixed};
\draw[ed] (b2.south) -- (rep.north);

% =============================================================== level 3 ==
\node[hd] at (0.4,-7.20) {\textbf{level 3} --- one delivery, once every
  $n_{gt}$ is fixed: the one furthest from an integer};

\node[nd, minimum width=32mm] (a3) at (4.9,-8.65)
  {$\overline{V}_{nt}\le\lfloor v\rfloor$};
\node[nd, minimum width=32mm] (c3) at (10.1,-8.65)
  {$\overline{V}_{nt}\ge\lceil v\rceil$};
\coordinate (j3) at (7.5,-8.05);
\draw[sp] (rep.south) -- (j3);
\draw[ed] (j3) -- (a3.north);
\draw[ed] (j3) -- (c3.north);

\node[an, text width=104mm] at (7.5,-9.25) {a delivery counts physical RTIs,
  so the two children are disjoint and cover $\mathbb{Z}$ between them; this
  rung is entered only where the first two have nothing left to say};

% --------------------------------------------------------- the non-rung --
\node[bnx] (ban2) at (7.5,-10.19) {\textbf{Not a level:} a single column
  variable $\lambda^{g}_{\omega t}$. A bound on one is not a row of the
  master --- honouring it means excluding that column from the pricing
  problem --- and the relaxation answers by moving the same weight to a
  column that differs in one RTI: $3\,000$ branches, no integral relaxation
  reached.};

\end{tikzpicture}
\caption{Every rung bounds an integer quantity the master already prices,
so the children of a split are disjoint and together exhaustive and the search
gives up nothing by taking them. The order is forced by the objective:
transport costs $cK\cdot X$ and nothing else in the model counts trucks, so
the horizon total is the one aggregate that cannot be answered by moving
freight between periods for one period's holding cost, as a single dispatch
count or a single delivery can. And because these are rows of the master and
not restrictions on the subproblem, branching moves the duals and leaves the
pricing problem alone --- which is also how a load no dual vector would have
asked for (Figure~\ref{fig:hull}) comes to be priced at all.}
\label{fig:tree}
\end{figure}
